\documentclass[11pt]{article}

\usepackage{amsmath,amsfonts}
\usepackage{graphicx}
\usepackage{hyperref}
\usepackage{geometry}
\geometry{margin=1in}

\title{Estimating Intrinsic Dimension of High-Dimensional Data:\\
A Comparative Study of Modern Estimators}

\author{
Caleb Fikes \\
Emory University
\and
Jack Michael Solomon \\
Emory University
}

\date{}

\begin{document}

\maketitle

\section*{Overview and Motivation}

Many machine learning datasets exist in extremely high-dimensional ambient spaces. For example, each MNIST digit image contains $28\times 28 = 784$ pixel features. However, the \emph{manifold hypothesis} posits that real-world data lie near a much lower-dimensional manifold embedded in the ambient space.

Formally, suppose observations are
\[
x_i \in \mathbb{R}^D.
\]
The manifold hypothesis assumes the data distribution is concentrated near a smooth manifold
\[
M \subset \mathbb{R}^D
\]
with intrinsic dimension
\[
d = \dim(M), \qquad d \ll D.
\]

Intrinsic dimension plays an important role in dimensionality reduction, representation learning, and theoretical analyses of sample complexity. For example, learning complexity may depend primarily on the intrinsic dimension rather than the ambient dimension \cite{narayanan2010}. Despite its importance, estimating intrinsic dimension from finite samples remains difficult.

Many estimators have been proposed based on nearest-neighbor statistics, correlation integrals, angular concentration, and geometric separability. However, these estimators often exhibit strong sensitivity to neighborhood size and noise. The goal of this project is to empirically evaluate several intrinsic dimension estimators and study their robustness across synthetic manifolds and real datasets.

\section*{Data}

We will evaluate estimators on both synthetic manifolds and real data.

\subsection*{Synthetic Data}

Synthetic datasets allow evaluation against known ground-truth intrinsic dimension. 
We will generate samples from two families of manifolds:

\begin{itemize}
\item Spheres $S^d$
\item $d$-dimensional tori $T^d = (S^1)^d$
\end{itemize}

Both manifolds have intrinsic dimension $d$, but differ in global geometry and topology. 
Comparing estimator performance across these families allows us to study whether estimators respond purely to intrinsic dimension or are sensitive to geometric structure.

Samples from $S^d$ will be generated by drawing
\[
z \sim \mathcal{N}(0, I_{d+1})
\]
and normalizing
\[
x = \frac{z}{\|z\|},
\]
which yields points uniformly distributed on the sphere.

Samples from the torus $T^d$ will be generated by drawing
\[
\theta_1,\dots,\theta_d \sim \mathrm{Uniform}[0,2\pi]
\]
and mapping them to
\[
x = (\cos\theta_1,\sin\theta_1,\dots,\cos\theta_d,\sin\theta_d)
\in \mathbb{R}^{2d}.
\]

To control the ambient dimension, all manifolds will be embedded into a fixed ambient space $\mathbb{R}^D$, where $D$ is chosen to match the largest natural embedding dimension in the experiment (for spheres $D = d_\text{max} +1$, and for tori $D=2d_{\max}$). Lower-dimensional manifolds will be embedded via random linear maps. Specifically, for each experiment we will draw a random orthonormal matrix and apply it to the sampled points to produce a random embedding.

Experiments will sweep over intrinsic dimensions $d$ while keeping the ambient dimension $D$ fixed. For each manifold and intrinsic dimension, we will repeat the experiment multiple times with independently sampled data and independently sampled random embeddings in order to measure estimator variability.

\subsection*{Real Data}

We will use the MNIST dataset consisting of 70,000 handwritten digit images of size $28\times28$, corresponding to ambient dimension $D=784$. The dataset therefore contains over $5\times10^7$ total feature values.

To study intrinsic dimension in learned representations, we will train autoencoders with varying bottleneck widths $k$. The encoder maps the data to a $k$-dimensional latent representation and the decoder reconstructs the input image.

If an autoencoder achieves near-zero reconstruction error with bottleneck dimension $k$, then the data can be represented using at most $k$ latent parameters. This provides an approximate upper bound on the intrinsic dimension of the learned representation. Consequently, any intrinsic-dimension estimate $\hat d > k$ indicates estimator inconsistency.

By varying bottleneck width and applying intrinsic-dimension estimators to the resulting latent representations, we can evaluate how estimator behavior changes when the representation dimension is controlled.

% \section*{Intrinsic Dimension Estimators}

% We will evaluate several estimators from the literature.

% \textbf{Levina–Bickel Maximum Likelihood Estimator} \cite{levina2005}

% Let $T_j(x_i)$ denote the distance from point $x_i$ to its $j$-th nearest neighbor. The local estimate is

% \[
% \hat d_i =
% \left(
% \frac{1}{k-1}
% \sum_{j=1}^{k-1}
% \log \frac{T_k(x_i)}{T_j(x_i)}
% \right)^{-1}.
% \]

% The global estimate is obtained by averaging over data points.

% \textbf{TwoNN Estimator} \cite{facco2017}

% Using the ratio of the first two nearest-neighbor distances

% \[
% \mu_i = \frac{T_2(x_i)}{T_1(x_i)},
% \]

% the intrinsic dimension estimate is

% \[
% \hat d = \left( \frac{1}{n} \sum_i \log \mu_i \right)^{-1}.
% \]

% \textbf{Correlation Dimension} \cite{grassberger1983}

% Define the correlation integral

% \[
% C(r) =
% \frac{2}{n(n-1)}\sum_{i<j}\mathbf{1}\{\|x_i-x_j\| < r\}.
% \]

% The intrinsic dimension is estimated from

% \[
% d = \frac{d \log C(r)}{d \log r}.
% \]

% \textbf{DANCo Estimator} \cite{ceruti2014}

% DANCo estimates intrinsic dimension by combining nearest-neighbor distance statistics with angular concentration in high dimensions.

% \textbf{MiND Estimator} \cite{rozza2012}

% MiND estimates intrinsic dimension using nearest-neighbor graph structure combined with manifold geometry assumptions.

% \textbf{Fisher Separability Dimension} \cite{gorban2019}

% This method estimates intrinsic dimension based on the separability of points by linear discriminants in high-dimensional space.

% \textbf{Neural Network Representation-Based Estimation}

% We will also study intrinsic dimension estimates on representations learned by autoencoders.

\section*{Intrinsic Dimension Estimators}

We evaluate several intrinsic-dimension estimators drawn from the literature. 
These estimators rely on different geometric properties of high-dimensional data, including nearest-neighbor distance statistics, fractal scaling laws, angular concentration phenomena, and high-dimensional linear separability.

\textbf{Levina--Bickel Maximum Likelihood Estimator} \cite{levina2005}

Let $T_j(x_i)$ denote the Euclidean distance from point $x_i$ to its $j$-th nearest neighbor. 
Under a local uniform sampling model on a $d$-dimensional manifold, the ratios of neighbor distances follow a distribution parameterized by $d$. 
The local intrinsic-dimension estimate is

\[
\hat d_i =
\left(
\frac{1}{k-1}
\sum_{j=1}^{k-1}
\log \frac{T_k(x_i)}{T_j(x_i)}
\right)^{-1}.
\]

The global estimate is obtained by averaging over all points:

\[
\hat d = \frac{1}{n}\sum_{i=1}^n \hat d_i .
\]

\textbf{TwoNN Estimator} \cite{facco2017}

Let $T_1(x_i)$ and $T_2(x_i)$ denote the distances from $x_i$ to its first and second nearest neighbors. Define

\[
\mu_i = \frac{T_2(x_i)}{T_1(x_i)}.
\]

Under a $d$-dimensional manifold model the random variable $\log \mu_i$ has expectation $1/d$. 
The intrinsic dimension estimate is therefore

\[
\hat d =
\left(
\frac{1}{n}\sum_{i=1}^n \log \mu_i
\right)^{-1}.
\]

\textbf{Correlation Dimension (CorrInt)} \cite{grassberger1983}

Define the correlation integral

\[
C(r) =
\frac{2}{n(n-1)}
\sum_{i<j}
\mathbf{1}\{\|x_i-x_j\| < r\}.
\]

For small $r$, the correlation integral scales as

\[
C(r) \propto r^d .
\]

Thus the intrinsic dimension can be estimated from the slope of the log--log scaling relation:

\[
\hat d =
\frac{d \log C(r)}{d \log r}.
\]

In practice this derivative is approximated by linear regression over a range of small radii.

\textbf{DANCo Estimator} \cite{ceruti2014}

DANCo estimates intrinsic dimension by jointly modeling nearest-neighbor distance statistics and angular concentration effects.

Let $\mathcal{R}$ denote the empirical distribution of normalized nearest-neighbor distances and $\mathcal{A}$ the empirical distribution of angles between neighbor difference vectors. 
For a candidate intrinsic dimension $d$, theoretical distributions $\mathcal{R}_d$ and $\mathcal{A}_d$ can be derived for points sampled from a $d$-dimensional space.

The intrinsic dimension estimate is obtained by minimizing the divergence between empirical and theoretical statistics:

\[
\hat d =
\arg\min_d
\left[
D_{\mathrm{KL}}(\mathcal{R}\,\|\,\mathcal{R}_d)
+
D_{\mathrm{KL}}(\mathcal{A}\,\|\,\mathcal{A}_d)
\right],
\]

where $D_{\mathrm{KL}}$ denotes the Kullback--Leibler divergence.

\textbf{MiND Estimator} \cite{rozza2012}

MiND estimates intrinsic dimension using nearest-neighbor distance statistics derived from the geometry of local neighborhoods.

Let $T_k(x_i)$ denote the distance from $x_i$ to its $k$-th nearest neighbor. 
Under a model in which points are sampled from a $d$-dimensional manifold, the distribution of normalized neighbor distances has a known parametric form depending on $d$.

The intrinsic dimension estimate is obtained by maximizing the likelihood of the observed neighbor distances:

\[
\hat d =
\arg\max_d
\sum_{i=1}^n
\log p\!\left(
T_1(x_i),\dots,T_K(x_i) \mid d
\right),
\]

where $p(\cdot \mid d)$ denotes the model likelihood under intrinsic dimension $d$.

\textbf{Fisher Separability Dimension} \cite{gorban2019}

This estimator is based on the probability that a point can be separated from the rest of the dataset by a linear discriminant.

Let $p_\alpha$ denote the empirical probability that a point $x_i$ cannot be linearly separated from the remaining points using a Fisher discriminant with margin parameter $\alpha$. 
Under a model where points lie on a $d$-dimensional sphere, the expected unseparability probability depends on $d$.

The intrinsic dimension estimate $\hat d$ is obtained by solving

\[
p_\alpha = p_d(\alpha),
\]

where $p_d(\alpha)$ denotes the theoretical unseparability probability for points sampled from a $d$-dimensional distribution.

\section*{Evaluation Methodology}

\textbf{Synthetic Data}

For synthetic manifolds with known intrinsic dimension we evaluate each estimator using Monte Carlo experiments. 
For a manifold family $M \in \{S^d, T^d\}$, intrinsic dimension $d$, neighborhood parameter $k$, noise level $\sigma$, and Monte Carlo replicate $r$, let

\[
\hat d_{E}(M,d,k,\sigma,r)
\]

denote the intrinsic-dimension estimate produced by estimator $E$.

Each experiment consists of independently sampling points on the manifold, embedding them into a fixed ambient space $\mathbb{R}^D$ via a random linear embedding, and optionally adding Gaussian noise orthogonal to the manifold. 
All reported quantities are computed over $R$ Monte Carlo repetitions.

\emph{Estimation Error.}  
We measure absolute estimation error using the mean absolute error

\[
\mathrm{MAE}_E(M,d,k,\sigma)
=
\frac{1}{R}
\sum_{r=1}^{R}
\left|
\hat d_{E}(M,d,k,\sigma,r)-d
\right|.
\]

\emph{Bias and Variance.}  
Let

\[
\bar d_E(M,d,k,\sigma)
=
\frac{1}{R}
\sum_{r=1}^{R}
\hat d_{E}(M,d,k,\sigma,r)
\]

denote the average estimate. 
The empirical bias and variance are

\[
\mathrm{Bias}_E(M,d,k,\sigma)
=
\bar d_E(M,d,k,\sigma)-d
\]

\[
\mathrm{Var}_E(M,d,k,\sigma)
=
\frac{1}{R-1}
\sum_{r=1}^{R}
\left(
\hat d_{E}(M,d,k,\sigma,r)
-
\bar d_E(M,d,k,\sigma)
\right)^2 .
\]

\emph{Neighborhood Sensitivity.}  
To evaluate stability with respect to the neighborhood parameter, we compute

\[
\mathrm{Stab}_E(M,d,\sigma)
=
\max_{k \in K}
\bar d_E(M,d,k,\sigma)
-
\min_{k \in K}
\bar d_E(M,d,k,\sigma),
\]

where $K$ is a grid of neighborhood sizes.

\emph{Robustness to Noise.}  
To measure sensitivity to noise we compute the change in the estimated dimension relative to the noiseless case:

\[
\mathrm{NoiseShift}_E(M,d,k,\sigma)
=
\bar d_E(M,d,k,\sigma)
-
\bar d_E(M,d,k,0).
\]

We perform dimension sweeps over $d$ for both spheres $S^d$ and tori $T^d$, repeating experiments across independently sampled datasets and independently drawn random embeddings in order to quantify estimator bias, variance, stability, and noise robustness.

\textbf{Real Data}

For real datasets we perform an analogous Monte Carlo evaluation using the MNIST dataset. 
Let $X \subset \mathbb{R}^{784}$ denote the set of MNIST images. 
To study intrinsic dimension under learned representations, we train autoencoders with bottleneck width $k$ and obtain latent representations

\[
z_i = f_k(x_i), \qquad z_i \in \mathbb{R}^k,
\]

where $f_k$ denotes the encoder of an autoencoder with bottleneck dimension $k$.

For estimator $E$, neighborhood parameter $k_n$, and Monte Carlo replicate $r$, let

\[
\hat d_E(k,k_n,r)
\]

denote the intrinsic-dimension estimate computed from the latent dataset $\{z_i\}$.

Each Monte Carlo replicate consists of independently training the autoencoder (with different random initialization and minibatch order) and computing the intrinsic-dimension estimates on the resulting latent representation.

\emph{Bottleneck Consistency Test.}  
If an autoencoder with bottleneck dimension $k$ achieves near-zero reconstruction error, then the dataset admits a learned representation whose effective dimensionality is at most $k$. 
Consequently, any estimate satisfying

\[
\hat d_E(k,k_n,r) > k
\]

indicates estimator inconsistency. 
We quantify this failure mode using the empirical violation rate

\[
\mathrm{ViolationRate}_E(k,k_n)
=
\frac{1}{R}
\sum_{r=1}^{R}
\mathbf{1}\{\hat d_E(k,k_n,r) > k\}.
\]

\emph{Estimator Variability.}  
To measure estimator stability across training runs we compute

\[
\bar d_E(k,k_n)
=
\frac{1}{R}
\sum_{r=1}^{R}
\hat d_E(k,k_n,r),
\]

\[
\mathrm{Var}_E(k,k_n)
=
\frac{1}{R-1}
\sum_{r=1}^{R}
\left(
\hat d_E(k,k_n,r)
-
\bar d_E(k,k_n)
\right)^2 .
\]

\emph{Neighborhood Sensitivity.}  
As in the synthetic experiments, we evaluate sensitivity to the neighborhood parameter using

\[
\mathrm{Stab}_E(k)
=
\max_{k_n \in K}
\bar d_E(k,k_n)
-
\min_{k_n \in K}
\bar d_E(k,k_n),
\]

where $K$ denotes a grid of neighborhood sizes.

\emph{Sample-Size Sensitivity.}
To measure robustness to finite-sample effects, we compute intrinsic-dimension estimates on random subsamples of the latent dataset of varying size $n$. Let
\[
\bar d_E(k,k_n,n)
=
\frac{1}{R}\sum_{r=1}^R \hat d_E(k,k_n,n,r).
\]
We then define
\[
\mathrm{SampleStab}_E(k,k_n)
=
\max_{n\in N}\bar d_E(k,k_n,n)
-
\min_{n\in N}\bar d_E(k,k_n,n),
\]
where $N$ is a grid of subsample sizes.

By sweeping over bottleneck widths $k$ and repeating experiments across independently trained autoencoders, we analyze estimator consistency, variability, and stability on real-world data representations.

% \section*{Proposed Contribution}

% This project will provide a systematic empirical comparison of intrinsic dimension estimators across synthetic manifolds and real datasets. We will analyze estimator stability with respect to neighborhood size and noise. If time permits, we will attempt to develop a stability-based aggregation method that produces more robust intrinsic dimension estimates.

\section*{Proposed Contribution}

This project will provide a systematic empirical comparison of intrinsic dimension estimators across synthetic manifolds and real datasets. Using Monte Carlo experiments on spheres $S^d$ and tori $T^d$ embedded into a fixed ambient space, we will evaluate estimator performance across a range of intrinsic dimensions while controlling for sampling variability, noise, and random embeddings. We will quantify estimator behavior using bias, variance, neighborhood stability, and noise sensitivity statistics, providing a controlled setting in which the reliability of existing estimators can be directly compared.

On real datasets, we will evaluate these estimators on learned representations of MNIST produced by autoencoders with varying bottleneck widths. By sweeping over bottleneck dimensions and repeatedly training autoencoders with different random initializations, we obtain multiple latent representations whose effective dimensionality is constrained by the bottleneck size. This enables a consistency diagnostic: if a representation admits near-zero reconstruction error with bottleneck dimension $k$, then any intrinsic-dimension estimate $\hat d > k$ indicates estimator inconsistency. In addition to this bottleneck test, we will measure estimator stability across neighborhood sizes and across random subsamples of the latent dataset in order to assess robustness to finite-sample effects.

Beyond benchmarking existing methods, we aim to investigate the structural properties that make certain estimators more stable than others. In particular, we will analyze how estimator behavior changes with neighborhood size, noise level, and sample size. Guided by these observations, we aim to develop and propose a stability-based aggregation strategy that combines estimates across neighborhood scales to produce more robust intrinsic-dimension estimates. Even if a new estimator is not ultimately proposed, our experiments will provide insight into which existing methods perform most reliably and under what conditions.
\bibliographystyle{plain}
\bibliography{references}

\end{document}